\scanpage{43}
\mspage{9}
elementary of $S^2$ \inserted{of (23)}, one therefore finds
that the datum of a gerbe on $S^2$ is equivalent
to that of an \uline{automorphism} of
the trivial extension \deleted{\unsure{over}} of $\mathbb Z$ by $A$,
hence \unsure{to|:} an element of $\operatorname{Hom}(\unsure{\mathbb Z|Z},A)\simeq A$,
whence the correct value
\[
(25)\qquad H^2(S^2,\deleted{\unsure{\mathbb Z|Z}}\inserted{A})\simeq A,
\]
magnificent success! Without going any
further in verifications, this
convinces me that the intuition that there
is a theory of \deleted{\unsure{inductive limits}\unsure{,}}\inserted{$\underset{\longrightarrow}{\mathrm{Lim}}$}
\unsure{of|of} homotopy types (\unsure{with|or} a
theorem \unsure{in|:} the \uline{van Kampen} style, saying
that the homotopy type of a
$\underline{\mathrm{Lim}}$ of Toposes (``glued'' Topos) is
the $\underline{\mathrm{Lim}}$ of the homotopy types),
is indeed correct. In principle, for
\margin{42}
